\scanpage{271}\mspage{\unsure{187}}
% archive sheet number: 270
4.7.20. Une \unsure{fcg|feg} canonique est dite \uline{maximale}
\inserted{(ou \uline{strictement canonique})},
si elle n'est pas \deleted{\unsure{déduit}} déduite d'une fcg canonique
str. plus grande $(\mathcal F^\pi)_{\pi\geq\pi_{00}}$ (donc avec
$\pi_{00}<\pi_0$). \deleted{plus} Donc si $M$ admet une
fcga, alors elle admet même une
fcg str. canonique, et cette dernière est de
plus unique. \deleted{\illegible{3}} On va revenir
là-dessus dans la section suivante.

4.8. \uline{Catégories} $\unsure{U|\mathcal U}$-\uline{accessibles}.

\deleted{\illegible{1}} \inserted{\uline{Proposition}} 4.8.1. Soit $M$ une $U$-catégorie,
\inserted{$\pi_0$ un cardinal infini. Considérons}
\deleted{\unsure{a)}} les conditions suivantes : \deleted{sont équivalentes :}

\deleted{\unsure{b)}}a) $M$ satisfait \unsure{à|:} $L\pi_0$ (i.e. est stable par \inserted{petites} $\varinjlim$
\unsure{grdes} devant $\pi_0$), et contient une petite
sous-catégorie pleine $\unsure{S|\mathcal S}$ \uline{strictement
génératrice}, formée d'objets $\pi_0$-accessibles.

b\deleted{\unsure{c}}) Il existe une petite catégorie $C$,
et une \deleted{\unsure{foncteur} pl. fidèle \unsure{non accessible|non essentiel}} \inserted{équivalence de catégories}
\[
M\approx\operatorname{Ind}_{\pi_0}(C)
\]

\deleted{\unsure{b)}}c) $M$ admet une filtration cardinale
généralisée accessible $(\mathcal F^\pi)_{\pi\geq\pi_0}$, ayant
$\pi_0$ comme cardinal initial, et telle que
$\pi_0$ soit un cardinal utile pour ladite
filtration.
